\chapter{Phosphate (Urine)}

\section*{What Is This Test?}
The urine phosphate test measures the amount of phosphate excreted in the urine, typically as a 24-hour collection or as a fractional excretion calculation from a random spot sample. Urinary phosphate excretion is the primary regulator of serum phosphate homeostasis. Approximately 80--90\% of filtered phosphate is reabsorbed in the proximal tubule via sodium-phosphate cotransporters (NaPi-IIa and NaPi-IIc), which are regulated by parathyroid hormone (PTH) and fibroblast growth factor 23 (FGF23) \cite{bergwitz2010regulation}. PTH decreases proximal tubular phosphate reabsorption by causing internalization and degradation of NaPi-IIa cotransporters, leading to phosphaturia. FGF23, produced by osteocytes, is the most potent phosphaturic hormone: it inhibits NaPi-IIa expression and suppresses 1,25-dihydroxyvitamin D (calcitriol) synthesis. The remaining filtered phosphate is reabsorbed in the distal nephron. The fractional excretion of phosphate (FEPO4) is normally 5--20\%, meaning that under normal conditions, the kidney excretes 5--20\% of filtered phosphate while reabsorbing 80--95\%. In hypophosphatemia, FEPO4 should decrease to <5\% as the kidney conserves phosphate. An FEPO4 >5\% in a hypophosphatemic patient indicates inappropriate renal phosphate wasting \cite{wallach2022interpretation}. Urine phosphate measurement is essential for evaluating phosphate-wasting disorders, parathyroid disorders, and the differential diagnosis of hypophosphatemia.

\section*{Alternative Names}
Urinary Phosphate; Urine PO4; 24-Hour Urine Phosphorus; Fractional Excretion of Phosphate (FEPO4); Tubular Reabsorption of Phosphate (TRP); Renal Phosphate Excretion; Spot Urine Phosphate/Creatinine Ratio

\section*{Principle}
Urine phosphate is measured by the same phosphomolybdate spectrophotometric method used for serum phosphate: phosphate reacts with ammonium molybdate in acidic solution to form phosphomolybdate, which is reduced to molybdenum blue (absorbance at 660--700 nm, proportional to phosphate concentration) \cite{tietz2012textbook}. The urine sample must be acidified to dissolve any precipitated calcium phosphate before analysis. For 24-hour collections, total phosphate is calculated from concentration $\times$ total volume. For spot samples, the fractional excretion of phosphate (FEPO4) is calculated: FEPO4 = (urine PO4 $\times$ serum creatinine) / (serum PO4 $\times$ urine creatinine) $\times$ 100\%. The tubular reabsorption of phosphate (TRP) is: TRP = 1 -- FEPO4. TRP >80\% is normal. FEPO4 calculation requires simultaneous blood and urine samples collected within a close timeframe. The test is performed on automated chemistry analyzers using the same methods as serum phosphate with appropriate urine sample dilution.

\section*{History}
The role of the kidney in phosphate homeostasis was established by Fuller Albright in the 1940s through studies of hypophosphatemic rickets. The Fiske-SubbaRow method for phosphate measurement was developed in 1925. The concept of fractional excretion and tubular maximum for phosphate (TmP/GFR) was developed in the 1960s. FGF23 was discovered in 2000 by the ADHR Consortium through genetic studies of autosomal dominant hypophosphatemic rickets, revolutionizing the understanding of phosphate disorders.

\section*{Why This Test Is Ordered}
\begin{itemize}
  \item Evaluation of hypophosphatemia: determine whether hypophosphatemia is due to renal phosphate wasting (FEPO4 >5--20\% inappropriately high for the level of hypophosphatemia) versus GI loss or transcellular shift (FEPO4 <5\% appropriately low)
  \item Diagnosis of X-linked hypophosphatemic rickets (XLH): FGF23 excess from PHEX mutation causes severe renal phosphate wasting, hypophosphatemia, low calcitriol, rickets, and osteomalacia. FEPO4 is elevated despite hypophosphatemia \cite{carpenter2011clinicians}
  \item Evaluation of oncogenic osteomalacia (tumor-induced osteomalacia): rare FGF23-secreting mesenchymal tumors cause acquired renal phosphate wasting, severe hypophosphatemia, bone pain, fractures, and muscle weakness \cite{chong2011tumor}
  \item Assessment of parathyroid disorders: primary hyperparathyroidism causes renal phosphate wasting (elevated FEPO4) with hypophosphatemia; hypoparathyroidism causes increased phosphate reabsorption (low FEPO4) with hyperphosphatemia
  \item Monitoring of chronic kidney disease: FEPO4 increases progressively as GFR declines to maintain phosphate balance; in CKD stage 4--5, FEPO4 often exceeds 50--70\%
  \item Evaluation of Fanconi syndrome: generalized proximal tubular dysfunction including phosphaturia, aminoaciduria, glycosuria, and bicarbonate wasting
\end{itemize}

\section*{Primary vs Secondary Test}
Urine phosphate measurement is a secondary, specialized test ordered when serum phosphate abnormalities require further characterization, particularly to determine whether renal phosphate wasting is present. It is not a routine test.

\section*{How This Test Is Performed}
\textbf{24-hour urine collection:} The patient collects all urine over a 24-hour period. The collection container is kept refrigerated and must be acidified (HCl) to prevent precipitation of calcium phosphate. \textbf{Random spot urine for FEPO4:} A random urine sample is collected simultaneously with a blood sample for serum phosphate and creatinine. Both urine and serum phosphate and creatinine must be measured from specimens drawn within a short timeframe (ideally within 1 hour). FEPO4 = (U\textsubscript{PO4} $\times$ P\textsubscript{Cr}) / (P\textsubscript{PO4} $\times$ U\textsubscript{Cr}) $\times$ 100\%, where U = urine concentration and P = plasma/serum concentration. The fractional excretion calculation from a spot sample is more convenient than a 24-hour collection and is accurate for assessing renal phosphate handling.

\section*{What Preparation Is Needed}
For an accurate FEPO4 or 24-hour urine phosphate, the patient should be on their usual diet. A very-low-phosphate diet will reduce urine phosphate; a high-phosphate diet will increase it. Fasting is not required for the spot FEPO4. Inform the provider of medications affecting urine phosphate: (1) Increase phosphaturia: PTH, loop diuretics, acetazolamide, glucocorticoids, calcitonin, bisphosphonates, phosphate binders, cinacalcet, denosumab. (2) Decrease phosphaturia (increase reabsorption): vitamin D/calcitriol, growth hormone, thyroid hormone, thiazide diuretics, bisphosphonates (variable effect). Time of day: urine phosphate excretion has a diurnal rhythm (highest in the afternoon, lowest in early morning) \cite{portale1987dietary}. Ideally, samples should be collected at the same time of day for repeat measurements.

\section*{Contraindications and Precautions}
No absolute contraindications. Important considerations: (1) FEPO4 values must be interpreted in the clinical context of the serum phosphate level. An FEPO4 of 10\% may be normal if serum phosphate is normal, but it is inappropriately high if the patient is hypophosphatemic (PO4 <2.5 mg/dL), where FEPO4 should be <5\%. (2) The tubular maximum reabsorption of phosphate per GFR (TmP/GFR = TRP $\times$ serum PO4) is a more precise index of renal phosphate handling than FEPO4. TmP/GFR <2.5 mg/dL indicates renal phosphate wasting. Nomograms are available \cite{walton1975nomogram}. (3) Incomplete 24-hour collection invalidates results---check 24-hour urine creatinine against expected range. (4) FEPO4 is influenced by dietary phosphate intake in the preceding days. (5) Acidified urine container is essential---if not acidified, calcium phosphate precipitates, falsely lowering measured phosphate. (6) FEPO4 changes with age: higher in children (growth, higher PTH), lower in elderly.

\section*{Risks}
No direct risk from urine collection. A blood draw is required for simultaneous serum phosphate and creatinine measurements for FEPO4 calculation; risks are minimal (mild pain/bruising from venipuncture).

\section*{What Happens After the Test}
Results available in 1--2 days. Interpretation of FEPO4 in hypophosphatemia: FEPO4 <5\% in a hypophosphatemic patient indicates appropriate renal phosphate conservation, suggesting that the cause of hypophosphatemia is GI loss (diarrhea, phosphate binders, vitamin D deficiency with decreased absorption), transcellular shift (refeeding syndrome, insulin therapy, respiratory alkalosis, hungry bone syndrome), or decreased intake. FEPO4 >5--20\% in a hypophosphatemic patient indicates inappropriate renal phosphate wasting, suggesting primary hyperparathyroidism, FGF23-mediated phosphate wasting (XLH, oncogenic osteomalacia), Fanconi syndrome, or post-renal transplant \cite{imel2012approach}. TmP/GFR is often reported and divides results into three groups: (a) NL (normal, >3.5 mg/dL), appropriate renal phosphate conservation; (b) LOW (<2.5 mg/dL), renal phosphate wasting; (c) marginal (2.5--3.5 mg/dL). For XLH and oncogenic osteomalacia, TmP/GFR is typically <2.0--2.5 mg/dL with hypophosphatemia and low or inappropriately normal calcitriol.

\section*{Understanding Results}

\subsection*{Below Normal}
\begin{itemize}
  \item \textbf{Low urine phosphate / FEPO4 <5\%:} Appropriate renal phosphate conservation.
  \item \textbf{GI phosphate loss or decreased absorption:} Chronic diarrhea, phosphate binders (calcium carbonate, aluminum hydroxide, sevelamer, lanthanum), vitamin D deficiency, malabsorption syndromes.
  \item \textbf{Transcellular shift:} Insulin/glucose administration, refeeding syndrome, respiratory alkalosis, hungry bone syndrome, catecholamines.
  \item \textbf{Decreased dietary phosphate intake:} Severe malnutrition, alcoholism, TPN without adequate phosphate.
  \item \textbf{Hypoparathyroidism:} Low PTH increases renal phosphate reabsorption (FEPO4 very low at <1--3\%). Serum phosphate is elevated.
\end{itemize}

\subsection*{Above Normal}
\begin{itemize}
  \item \textbf{FEPO4 >5--20\% in hypophosphatemia:} Inappropriate renal phosphate wasting. FEPO4 >20\% with normal phosphate indicates significant phosphaturia.
  \item \textbf{Primary hyperparathyroidism:} Excess PTH downregulates NaPi-IIa, causing phosphaturia. Serum Ca elevated, PO4 low, PTH elevated, calcitriol normal or elevated. FEPO4 elevated (typically 15--30\%).
  \item \textbf{FGF23-mediated disorders:} (a) X-linked hypophosphatemic rickets (XLH): PHEX gene mutation $\rightarrow$ elevated FGF23 $\rightarrow$ severe phosphaturia, hypophosphatemia, low calcitriol, normal calcium, normal/low PTH, rickets in children, osteomalacia in adults. FEPO4 markedly elevated. (b) Autosomal dominant hypophosphatemic rickets (ADHR): mutation in FGF23 cleavage site preventing its degradation. (c) Oncogenic osteomalacia: FGF23-secreting mesenchymal tumors (usually benign phosphaturic mesenchymal tumors) causing acquired XLH-like picture in adults---severe bone pain, muscle weakness, fractures, hypophosphatemia, low calcitriol, normal calcium, extremely elevated FEPO4. Locating the small tumor is challenging---often requires whole-body MRI, octreotide scan, or PET-CT. Surgical resection is curative.
  \item \textbf{Fanconi syndrome:} Generalized proximal tubular dysfunction from: cystinosis (most common inherited cause in children), Wilson disease, heavy metals (lead, cadmium, mercury), tenofovir disoproxil fumarate (anti-HIV medication---fanconi syndrome in 1--2\% of patients), ifosfamide, outdated tetracycline, multiple myeloma (light chain proximal tubular toxicity), Sj\"{o}gren's syndrome, Lowe syndrome, Dent disease. FEPO4 elevated along with glycosuria, aminoaciduria, uricosuria, and bicarbonate wasting (proximal RTA).
  \item \textbf{Post-renal transplant:} Persistent hyperparathyroidism from CKD plus glucocorticoid-induced phosphaturia.
  \item \textbf{Hypercalcemia of any cause:} Calcium competes at the calcium-sensing receptor in the proximal tubule, decreasing phosphate reabsorption.
  \item \textbf{Acetazolamide therapy:} Carbonic anhydrase inhibitor causing proximal tubular phosphate loss.
  \item \textbf{Chronic kidney disease:} FEPO4 rises progressively as functional nephron mass declines, reaching >50--70\% in ESRD. This is a compensatory mechanism to maintain phosphate balance, driven by rising PTH and FGF23.
\end{itemize}

\section*{Normal Results}
\begin{itemize}
  \item \textbf{24-hour urine phosphate (adults, normal diet):} 400--1,300 mg/24h (13--42 mmol/24h)---varies widely with dietary intake
  \item \textbf{Fractional excretion of phosphate (FEPO4, adults):} 5--20\%
  \item \textbf{FEPO4 (in hypophosphatemia---appropriate renal response):} <5\%
  \item \textbf{Tubular reabsorption of phosphate (TRP):} 80--95\%
  \item \textbf{TmP/GFR (tubular maximum reabsorption of phosphate):} 2.5--4.2 mg/dL (0.8--1.35 mmol/L)
  \item \textbf{Children:} FEPO4 10--30\% (higher due to growth-related phosphate demand and higher PTH status)
  \item \textbf{Urine phosphate/creatinine ratio (spot, adults):} <0.7 mg/mg
\end{itemize}

\section*{Follow-up Tests}
\begin{itemize}
  \item \textbf{Serum phosphate, calcium, PTH:} Essential for differential diagnosis of renal phosphate wasting.
  \item \textbf{25-hydroxyvitamin D and 1,25-dihydroxyvitamin D:} Low calcitriol with elevated FEPO4 suggests FGF23-mediated process.
  \item \textbf{Serum FGF23 (intact and C-terminal):} Elevated in XLH, ADHR, oncogenic osteomalacia, and CKD. Order if suspected FGF23-mediated phosphate wasting with no other explanation.
  \item \textbf{Serum creatinine and eGFR:} Required for accurate FEPO4 and TmP/GFR calculation.
  \item \textbf{Complete metabolic panel:} Potassium (often low in Fanconi syndrome with bicarbonate wasting), glucose, uric acid (low in Fanconi).
  \item \textbf{Genetic testing:} PHEX gene for XLH, FGF23 for ADHR, SLC34A3 for hereditary hypophosphatemic rickets with hypercalciuria.
  \item \textbf{Whole-body imaging:} To localize tumor in suspected oncogenic osteomalacia (MRI, octreotide scan, FDG-PET/CT with DOTATATE).
  \item \textbf{Renal ultrasound:} Assess for nephrocalcinosis.
  \item \textbf{Bone-specific alkaline phosphatase (BSAP):} Elevated in rickets/osteomalacia from phosphate wasting.
\end{itemize}

\section*{References}
\bibliographystyle{unsrtnat}
\bibliography{references}

\clearpage
\section*{Regulatory Status and Authorization Date}
Regulatory status applies to the specific assay, instrument, manufacturer, and intended use rather than to the test name alone. No single approval date can be assigned to this test category. For a commercial assay, verify the current product in the relevant regulator's database: the U.S. Food and Drug Administration (FDA) clearance, approval, or authorization; India's Central Drugs Standard Control Organisation (CDSCO) licence or permission; the European Union In Vitro Diagnostic Regulation (EU IVDR) CE certificate issued through the applicable conformity-assessment route; or the United Kingdom Medicines and Healthcare products Regulatory Agency (MHRA) registration and UKCA/CE status. Laboratory-developed tests may not have an individual FDA, CDSCO, EU IVDR, or MHRA approval date.
\clearpage
